\chapter{Espaços Vetoriais}\label{ch:b1:vspaces}

A álgebra linear começa aqui: os axiomas dos \omterm{def:b1:vspaces:def}{espaços vetoriais} isolam
o que $\R^2$, $\R^3$, os espaços de \omterm{def:b1:poly:def}{polinômios} e os espaços de funções
têm em comum --- pode-se somar e multiplicar por escalares. Dois
capítulos constroem a teoria (o \cref{ch:b1:findim} acrescenta a dimensão); a
linguagem que eles instalam --- \omterm{def:b1:vspaces:span}{espaço gerado}, \omterm{def:b1:vspaces:free}{família livre}, \omterm{def:b1:vspaces:free}{base},
\omterm{def:b1:vspaces:sum}{soma direta} --- é o pão de cada dia de todos os capítulos seguintes.
Ao longo do texto, $K$ designa $\R$ ou $\C$ (\emph{escalares}).

\section{Definição e exemplos}

\begin{definition}[Espaço vetorial]\label{def:b1:vspaces:def}
Um \emph{espaço vetorial sobre $K$}\index{espaço vetorial} é um
\omterm{def:b1:logic:sets}{conjunto} $E$ munido de uma adição que faz de $(E, +)$ um \omterm{def:b1:structures:group}{grupo
abeliano} (o zero escrito $0_E$ ou $0$), e de uma multiplicação por
escalares $K \times E \to E$ tal que, para todos $\lambda, \mu \in K$ e $x, y \in E$:
\[
\lambda(x + y) = \lambda x + \lambda y,\quad
(\lambda + \mu) x = \lambda x + \mu x,\quad
\lambda(\mu x) = (\lambda\mu) x,\quad
1\,x = x .
\]
Consequências: $0\,x = 0_E$, $\lambda\,0_E = 0_E$, $(-1)x = -x$, e
$\lambda x = 0_E \implies \lambda = 0$ ou $x = 0_E$ (multiplique por
$\lambda^{-1}$).
\end{definition}

\begin{proof}[Demonstração das consequências]
Para $0\,x = 0_E$: de $(0 + 0)x = 0x + 0x$ e $(0+0)x = 0x$,
cancele $0x$ no \omterm{def:b1:structures:group}{grupo} $(E, +)$. Para $\lambda\,0_E$: o mesmo
truque em $\lambda(0_E + 0_E)$. Para $(-1)x$: some $x$,
\[
x + (-1)x = 1\,x + (-1)x = \bigl(1 + (-1)\bigr)x = 0\,x = 0_E ,
\]
logo $(-1)x$ é o inverso aditivo de $x$. Por fim, se $\lambda x
= 0_E$ com $\lambda \neq 0$:
multiplique por $\lambda^{-1}$ (os escalares formam um \omterm{def:b1:structures:field}{corpo}) e use os dois
axiomas $\lambda^{-1}(\lambda x) = (\lambda^{-1}\lambda)x = 1x = x$
junto com $\lambda^{-1}0_E = 0_E$: $x = 0_E$. Por menores que sejam,
estas quatro regras são usadas silenciosamente em cada página que se
segue --- e a última é exatamente onde os \omterm{def:b1:structures:field}{corpos} são necessários:
sobre os escalares $\Z$, o ``espaço'' $\Z/2\Z$ a violaria
com $2\,x = 0$.
\end{proof}

\begin{example}\label{ex:b1:vspaces:examples}
$K^n$ (operações coordenada a coordenada); os \omterm{def:b1:poly:def}{polinômios} $K[X]$; as
funções $\mathcal{F}(A, K)$ de um \omterm{def:b1:logic:sets}{conjunto} qualquer $A$ em $K$ (operações
ponto a ponto) --- contendo as funções \omterm{def:b1:continuity:continuous}{contínuas}, as sequências
$\mathcal{F}(\N, \R)$, etc.; $\C$ como \omterm{def:b1:vspaces:def}{espaço vetorial} sobre $\R$. Em cada
caso os axiomas são herdados dos de $K$.
\end{example}

\begin{definition}[Subespaço]\label{def:b1:vspaces:subspace}
$F \subseteq E$ é um \emph{subespaço}\index{subespaço} quando $0_E \in
F$ e $F$ é
estável por adição e por multiplicação por escalares ---
equivalentemente:
\[
F \neq \emptyset
\qquad\text{e}\qquad
\forall x, y \in F,\ \forall \lambda \in K,\quad
x + \lambda y \in F .
\]
Um subespaço é ele próprio um \omterm{def:b1:vspaces:def}{espaço vetorial}. Toda interseção de
subespaços é um subespaço; uma união quase nunca é (mesma
demonstração do \cref{exo:b1:structures:6}).
\end{definition}

\begin{example}\label{ex:b1:vspaces:subspaces}
Em $\mathcal{F}(\R, \R)$: as funções \omterm{def:b1:continuity:continuous}{contínuas}, as deriváveis, os
\omterm{def:b1:poly:def}{polinômios} de grau $\leq n$ (escrito $K_n[X]$ dentro de $K[X]$), as
soluções de uma equação diferencial linear homogênea (o \cref{thm:b1:diffeq:homogeneous2} dizia
exatamente isso). Não exemplos: $\{f : f(0) = 1\}$ (sem o zero), grau exatamente
$n$ (não estável por adição).
\end{example}

\begin{example}[Subespaço ou não: quatro veredictos, argumentados]
\label{ex:b1:vspaces:subspaceornot}
No espaço das sequências reais:
\begin{itemize}
  \item $\{u : u \text{ limitada}\}$ \emph{é} um \omterm{def:b1:vspaces:subspace}{subespaço}: $0$ é
        limitada, e se $\abs{u_n} \leq M$, $\abs{v_n} \leq M'$,
        então $\abs{u_n + \lambda v_n} \leq M + \abs\lambda M'$.
  \item $\{u : u_n \to 1\}$ \emph{não} é: falta a sequência nula (e a soma de
        dois membros tende a $2$).
  \item $\{u : u \text{ monotone}\}$ \emph{não} é: $u_n = n$
        e $v_n = -n + (-1)^n$ são monótonas, e a sua soma
        $(-1)^n$ não é; a estabilidade por adição é o axioma
        que falha, embora o \omterm{def:b1:logic:sets}{conjunto} contenha $0$ e todos os
        múltiplos escalares dos seus membros.
  \item $\{u : u_{n+1} = u_n^2\}$ \emph{não} é: contém $0$,
        mas $2u$ escapa assim que $u$ é um membro não nulo
        ($2u_{n+1} \neq (2u_n)^2$ em geral) --- a não linearidade
        está no quadrado.
\end{itemize}
A ordem de trabalho é sempre a mesma: teste $0$ primeiro (o mais
barato), depois a estabilidade --- e, para refutar, um único par
explícito de contraexemplos vale mais do que qualquer dose de dúvida.
\end{example}

\section{Espaço gerado, somas, somas diretas}

\begin{definition}[Combinações lineares, espaço gerado]\label{def:b1:vspaces:span}
Uma \emph{combinação linear} da família $(x_1, \dots, x_p)$ de
vetores de $E$ é qualquer $\lambda_1 x_1 + \dots + \lambda_p x_p$
($\lambda_i \in K$). O \omterm{def:b1:logic:sets}{conjunto} de todas elas é o
\emph{espaço gerado}\index{espaço gerado} $\operatorname{Vect}(x_1, \dots, x_p)$: é
um \omterm{def:b1:vspaces:subspace}{subespaço}, o menor que contém a família.
\end{definition}

\begin{proof}[Demonstração das duas afirmações]
Estabilidade: uma soma de duas combinações lineares $\sum\lambda_i x_i +
\sum\mu_i x_i = \sum(\lambda_i + \mu_i)x_i$ é ainda
uma, e um múltiplo escalar $\mu\sum\lambda_i x_i = \sum(\mu\lambda_i)
x_i$ também; a combinação nula mostra que
$0$ pertence: o \omterm{def:b1:vspaces:span}{espaço gerado} é um \omterm{def:b1:vspaces:subspace}{subespaço}. Minimalidade: seja $H$
um \omterm{def:b1:vspaces:subspace}{subespaço} qualquer contendo $x_1,
\dots, x_p$. Pela estabilidade por
multiplicação por escalares, cada $\lambda_i x_i \in H$, e pela estabilidade por
adição a sua soma pertence a $H$: toda combinação linear pertence a
$H$, isto é, $\operatorname{Vect}(x_1, \dots, x_p) \subseteq H$. Assim o \omterm{def:b1:vspaces:span}{espaço gerado} está
contido em todo \omterm{def:b1:vspaces:subspace}{subespaço} que contém a família: é o menor deles.
\end{proof}

\begin{definition}[Soma, soma direta]\label{def:b1:vspaces:sum}
Para \omterm{def:b1:vspaces:subspace}{subespaços} $F, G$ de $E$:
\[
F + G = \{\,u + v : u \in F,\ v \in G\,\}
\]
é um \omterm{def:b1:vspaces:subspace}{subespaço} (o menor que contém $F \cup G$). A soma é
\emph{direta}\index{soma direta}, escrita $F \oplus G$, quando todo
elemento de $F + G$ se decompõe de modo \emph{único} como $u + v$;
equivalentemente (veja abaixo) quando $F \cap G = \{0\}$. Quando $E = F \oplus
G$, os
\omterm{def:b1:vspaces:subspace}{subespaços} são \emph{suplementares}\index{subespaços suplementares}
em $E$.
\end{definition}

\begin{example}[Uma soma de duas retas]
\label{ex:b1:vspaces:sumoflines}
Em $\R^3$, sejam $F = \operatorname{Vect}\bigl((1,0,1)\bigr)$ e
$G = \operatorname{Vect}\bigl((0,1,1)\bigr)$. A sua soma é
\[
F + G = \{\,a(1,0,1) + b(0,1,1)\,\}
= \{(a,\ b,\ a + b)\} = \{(x, y, z) : z = x + y\},
\]
o plano pela origem que contém as duas retas. Ele é estritamente
maior do que a união $F \cup G$ (a mera cruz das duas retas): o vetor
$(1, 1, 2) = (1,0,1) + (0,1,1)$ está na soma mas em nenhuma das duas retas. E $F \cap G = \{0\}$ (um vetor
comum exigiria $a(1,0,1) = b(0,1,1)$, cujas duas primeiras \omterm{prop:b1:vspaces:coordinates}{coordenadas} forçam $a = b = 0$):
a soma é direta, e $F \oplus
G$ é exatamente aquele plano.
\end{example}

\begin{proposition}\label{prop:b1:vspaces:directsum}
$F + G$ é direta se e somente se $F \cap G = \{0\}$.
\end{proposition}

\begin{proof}
Se algum $w \neq 0$ pertence a $F \cap G$: $w = w + 0 = 0 + w$ são duas
decomposições de $w$. Reciprocamente, se $u + v = u' + v'$ com $u, u'
\in F$, $v, v' \in G$,
então $u - u' = v' - v$ pertence a $F \cap G =
\{0\}$: as decomposições são únicas.
\end{proof}

\begin{method}[Demonstrar que $E = F \oplus G$]
\label{met:b1:vspaces:directsum}
Duas coisas a verificar, cada uma com o seu lance de abertura padrão.
\begin{enumerate}
  \item \emph{Interseção trivial.} Tome $x \in F \cap G$,
        escreva as duas condições de pertinência e esprema até
        $x =
        0$. (Nunca argumente por desenho: cf.\ as armadilhas
        abaixo.)
  \item \emph{A soma é tudo.} Tome um $x \in E$ arbitrário
        e \emph{produza} a decomposição $x = f + g$ ---
        seja adivinhando $f$ a partir do alvo ($f$ deve satisfazer
        a propriedade que define $F$, o que em geral dita a sua
        fórmula), seja resolvendo o sistema linear que exprime $x$
        contra geradores de $F$ e de $G$.
\end{enumerate}
Quando a fórmula da decomposição é adivinhada, a unicidade é
automática pelo passo 1; quando só a existência é duvidosa, é no
passo 2 que mora o trabalho. Os dois exemplos abaixo executam o
método: para as funções pares e ímpares, a fórmula de $f$ é forçada
avaliando a identidade desejada em $x$ e em $-x$; para os \omterm{def:b1:poly:def}{polinômios}
que se anulam num ponto, avaliando em $a$.
\end{method}

\begin{example}\label{ex:b1:vspaces:evenodd}
Em $\mathcal{F}(\R, \R)$, as funções pares $\mathcal{P}$ e as funções ímpares $\mathcal{I}$ são
\omterm{def:b1:vspaces:sum}{suplementares}: toda $f$ se escreve
\[
f(x) = \underbrace{\frac{f(x) + f(-x)}{2}}_{\text{par}}
+ \underbrace{\frac{f(x) - f(-x)}{2}}_{\text{ímpar}},
\]
e uma função ao mesmo tempo par e ímpar é nula. (Aplicado a $\exp$,
isto é o par $(\cosh, \sinh)$ do \cref{ch:b1:functions}.)
\end{example}

\begin{example}[{Um par suplementar em $K_n[X]$}]
\label{ex:b1:vspaces:hyperplane}
Fixe $a \in K$ e ponha $F = \{P \in K_n[X] : P(a) = 0\}$, $G =
\operatorname{Vect}(1)$ (as constantes). Então $K_n[X] = F \oplus
G$. De
fato $F \cap G$ consiste nas constantes que se anulam em $a$, isto é,
$\{0\}$; e todo $P$ se decompõe como
\[
P = \underbrace{\bigl(P - P(a)\bigr)}_{\in F}
+ \underbrace{P(a)}_{\in G} .
\]
Vale a pena memorizar a decomposição: subtrair o valor num ponto é a
maneira padrão de projetar sobre as ``funções que se anulam em
$a$''. Note que $F$ é um \omterm{def:b1:vspaces:subspace}{subespaço} grande e $G$ um pequeno; um par
\omterm{def:b1:vspaces:sum}{suplementar} não precisa ser equilibrado em sentido algum.
\end{example}

\begin{example}[Um subespaço suplementar nunca é único]
\label{ex:b1:vspaces:manysupplements}
Em $\R^2$, seja $F = \operatorname{Vect}\bigl((1,0)\bigr)$ (o eixo $x$). Tanto $G = \operatorname{Vect}\bigl((0,1)\bigr)$ quanto $G' = \operatorname{Vect}\bigl((1,1)\bigr)$ são
\omterm{def:b1:vspaces:sum}{suplementares} de $F$: cada um encontra $F$ apenas em $0$, e cada par
soma $\R^2$. As decomposições de um mesmo vetor diferem:
\[
(2,\ 1.5) = \underbrace{(2, 0)}_{\in F} +
\underbrace{(0, 1.5)}_{\in G}
= \underbrace{(0.5,\ 0)}_{\in F} +
\underbrace{(1.5,\ 1.5)}_{\in G'} .
\]
De fato, \emph{toda} reta distinta do próprio $F$ é um \omterm{def:b1:vspaces:sum}{suplementar} de
$F$ em $\R^2$: os \omterm{def:b1:vspaces:sum}{suplementares} são abundantes, e falar em ``o''
\omterm{def:b1:vspaces:sum}{suplementar} não faz sentido enquanto uma estrutura adicional (um
produto interno, \cref{ch:b1:euclid}) não singularizar um deles.
\end{example}

\begin{omfigure}
\begin{tikzpicture}[scale=1.5]
  \draw[->, thin] (-0.4,0) -- (2.9,0) node[below] {$x$};
  \draw[->, thin] (0,-0.4) -- (0,1.9) node[left] {$y$};
  \draw[omDef, very thick] (-0.3,0) -- (2.7,0)
    node[below, pos=0.95, font=\small] {$F$};
  \draw[omThm, very thick] (0,-0.3) -- (0,1.7)
    node[left, pos=0.95, font=\small] {$G$};
  \draw[omProp, very thick] (-0.25,-0.25) -- (1.75,1.75)
    node[right, pos=0.97, font=\small] {$G'$};
  \fill (2,1.5) circle (0.045) node[above right, font=\small]
    {$(2, 1.5)$};
  \draw[omThm, dashed] (2,1.5) -- (2,0);
  \draw[omDef, dashed] (2,1.5) -- (0,1.5);
  \draw[omProp, dashed] (2,1.5) -- (0.5,0);
  \fill (2,0) circle (0.035);
  \fill (0.5,0) circle (0.035);
\end{tikzpicture}

{\small Duas decomposições do mesmo ponto de $\R^2$ ao longo de $F$
(o eixo $x$): com o \omterm{def:b1:vspaces:sum}{suplementar} $G$ (queda vertical) e com o
\omterm{def:b1:vspaces:sum}{suplementar} $G'$ (queda oblíqua). As componentes em $F$ diferem:
uma projeção depende da direção da descida.}
\end{omfigure}

\section{Famílias livres, famílias geradoras, bases}

\begin{definition}\label{def:b1:vspaces:free}
Uma família $(x_1, \dots, x_p)$ de vetores de $E$ é:
\begin{itemize}
  \item \emph{geradora} (de $E$) quando
        $\operatorname{Vect}(x_1,\dots,x_p) = E$;
  \item \emph{livre}\index{família livre} (os seus vetores
        \emph{linearmente independentes}) quando
        \[
        \lambda_1 x_1 + \dots + \lambda_p x_p = 0
        \implies \lambda_1 = \dots = \lambda_p = 0 ;
        \]
        caso contrário, \emph{linearmente dependente};
  \item uma \emph{base}\index{base} quando é livre e geradora.
\end{itemize}
\end{definition}

\begin{proposition}[Coordenadas]\label{prop:b1:vspaces:coordinates}
$(e_1, \dots, e_n)$ é uma \omterm{def:b1:vspaces:free}{base} de $E$ se e somente se todo $x \in E$
é de modo \emph{único} uma combinação $x = \lambda_1 e_1 + \dots +
\lambda_n e_n$; os escalares $\lambda_i$ são
as \emph{coordenadas}\index{coordenadas} de $x$ na \omterm{def:b1:vspaces:free}{base}.
\end{proposition}

\begin{proof}
Geradora $=$ existência da decomposição. Unicidade $=$ liberdade:
duas decomposições do mesmo $x$ diferem por uma combinação igual a
$0$; a liberdade força todos os seus coeficientes --- as diferenças
de \omterm{prop:b1:vspaces:coordinates}{coordenadas} --- a se anularem. Reciprocamente, uma combinação nula
não trivial dá as duas decomposições $0 = \sum \lambda_i
e_i = \sum 0\,e_i$.
\end{proof}

\begin{example}\label{ex:b1:vspaces:bases}
A \emph{\omterm{def:b1:vspaces:free}{base} canônica} de $K^n$: $e_i = (0, \dots, 1, \dots, 0)$
($1$ na posição $i$). Os monômios $(1, X, X^2, \dots, X^n)$: uma \omterm{def:b1:vspaces:free}{base}
de $K_n[X]$ (liberdade: uma combinação nula é o \omterm{def:b1:poly:def}{polinômio} nulo, logo
todos os coeficientes se anulam, \cref{def:b1:poly:def}). Em
$\C$ sobre $\R$: a \omterm{def:b1:vspaces:free}{base} $(1, \iu)$.
\end{example}

\begin{remark}[As coordenadas são um trabalho de equipe]
\label{rem:b1:vspaces:teameffort}
A primeira coordenada de $x$ numa \omterm{def:b1:vspaces:free}{base} $(e_1, \dots, e_n)$
depende de \emph{todos} os vetores da \omterm{def:b1:vspaces:free}{base}, e não apenas de $e_1$. Em
$\R^2$: o vetor $(3, 1)$ tem primeira coordenada $3$ na \omterm{def:b1:vspaces:free}{base}
canônica, mas primeira coordenada $2$ na \omterm{def:b1:vspaces:free}{base}
$\bigl((1,0), (1,1)\bigr)$ --- resolva $(3,1) = a(1,0) + b(1,1)$:
$b = 1$, $a = 2$. Trocar um vetor da \omterm{def:b1:vspaces:free}{base} embaralha \emph{todas} as
\omterm{prop:b1:vspaces:coordinates}{coordenadas}; o \cref{ch:b1:matrices} empacotará esse embaralhamento na matriz de
mudança de \omterm{def:b1:vspaces:free}{base}.
\end{remark}

\begin{example}[Testar uma candidata a base, do início ao fim]
\label{ex:b1:vspaces:testbasis}
Será $\mathcal{F} = (1 + X,\ 1 + X^2,\ X + X^2)$ uma \omterm{def:b1:vspaces:free}{base} de
$\R_2[X]$? Escreva $u_1, u_2, u_3$ para os três \omterm{def:b1:poly:def}{polinômios}.
\emph{Liberdade}: uma combinação nula $a\,u_1 + b\,u_2 + c\,u_3 =
0$ dá, coeficiente a
coeficiente,
\[
a + b = 0, \qquad a + c = 0, \qquad b + c = 0 ;
\]
subtraindo as duas primeiras, $b = c$, e então a terceira dá $2b =
0$: $a = b = c = 0$, \omterm{def:b1:vspaces:free}{livre}. \emph{Geradora}: em vez de resolver três sistemas,
note a combinação simétrica
\[
u_1 + u_2 - u_3 = (1 + X) + (1 + X^2) - (X + X^2) = 2 ,
\]
logo $1 = \frac12(u_1 + u_2 - u_3)$; então
\[
X = u_1 - 1 = \tfrac12\bigl(u_1 - u_2 + u_3\bigr),
\qquad
X^2 = u_2 - 1 = \tfrac12\bigl(-u_1 + u_2 + u_3\bigr).
\]
Os monômios estão no \omterm{def:b1:vspaces:span}{espaço gerado}, logo tudo está: $\mathcal{F}$
é uma \omterm{def:b1:vspaces:free}{base}. Como bônus, juntar as três fórmulas exibidas dá as
\omterm{prop:b1:vspaces:coordinates}{coordenadas} de qualquer $P = \alpha + \beta X + \gamma X^2$:
\[
P = \frac{\alpha + \beta - \gamma}{2}\,u_1
+ \frac{\alpha - \beta + \gamma}{2}\,u_2
+ \frac{-\alpha + \beta + \gamma}{2}\,u_3 .
\]
(Verificação de coerência com $P = X$: \omterm{prop:b1:vspaces:coordinates}{coordenadas} $\bigl(\frac12,
-\frac12, \frac12\bigr)$, como
encontrado acima.) Duas lições: a simetria numa família costuma
esconder uma combinação atalho; e assim que a dimensão estiver
disponível (\cref{ch:b1:findim}), toda a metade ``geradora'' deste trabalho virá
sem custo algum --- três vetores \omterm{def:b1:vspaces:free}{livres} de um espaço de dimensão $3$
formam sempre uma \omterm{def:b1:vspaces:free}{base}.
\end{example}

\begin{proposition}[Critérios úteis de liberdade]\label{prop:b1:vspaces:freecriteria}
\begin{enumerate}
  \item Uma família de \emph{\omterm{def:b1:poly:def}{polinômios} não nulos de graus dois a
        dois distintos} é \omterm{def:b1:vspaces:free}{livre}.
  \item Acrescentar um vetor a uma \omterm{def:b1:vspaces:free}{família livre} mantém-na \omterm{def:b1:vspaces:free}{livre} se e
        somente se o vetor está fora do \omterm{def:b1:vspaces:span}{espaço gerado} pela família.
  \item Toda subfamília de uma \omterm{def:b1:vspaces:free}{família livre} é \omterm{def:b1:vspaces:free}{livre}; toda família
        que contém uma família geradora é geradora.
\end{enumerate}
\end{proposition}

\begin{proof}
(1) Numa combinação nula, olhe para o maior grau presente: o seu
coeficiente deve se anular (nada cancela esse grau), e desça em
cascata.

(2) Se $x \in \operatorname{Vect}(x_1, \dots, x_p)$, a relação $x
- \sum\lambda_i x_i = 0$ é não trivial. Reciprocamente, uma
combinação nula não trivial de $(x_1, \dots, x_p, x)$ deve envolver $x$ com
coeficiente não nulo (senão contradiz a liberdade da família
pequena), e isolar $x$ o coloca no \omterm{def:b1:vspaces:span}{espaço gerado}.

(3) Subfamília: uma combinação nula da subfamília é uma da família
inteira com os coeficientes ausentes postos a $0$; a liberdade da
família grande mata todos eles. Superfamília: todo vetor de $E$ já é
uma combinação da parte geradora; dê aos vetores extras o
coeficiente $0$.
\end{proof}

\begin{example}[O princípio da escada]
\label{ex:b1:vspaces:staircase}
Sejam $P_0, P_1, \dots, P_n \in K_n[X]$ com $\deg P_k = k$ para cada $k$ (uma
``escada'' de graus). Então $(P_0, \dots, P_n)$ é uma \omterm{def:b1:vspaces:free}{base} de $K_n[X]$. A liberdade
é a \cref{prop:b1:vspaces:freecriteria}~(1). Para a propriedade geradora, argumente por descida
finita sobre o grau: seja $Q \in
K_n[X]$, $Q \neq 0$, de grau $d$, com coeficiente
dominante $a$, e seja $b \neq 0$ o coeficiente dominante de $P_d$. Então
$Q -
\frac ab P_d$ tem grau $< d$ (os termos de topo cancelam-se); substituindo
$Q$ por essa diferença e iterando, após no máximo $n + 1$ passos
chega-se ao \omterm{def:b1:poly:def}{polinômio} nulo, e desfazer as subtrações exprime $Q$ como
combinação dos $P_k$. Duas escadas já encontradas: as potências
transladadas $\bigl((X-a)^k\bigr)_{0 \leq k
\leq n}$ (\cref{exo:b1:vspaces:4}), e os produtos de Newton
$\bigl((X - x_0)(X - x_1)\cdots(X - x_{k-1})\bigr)_{0 \leq k \leq
n}$, postos a trabalhar no problema de fim de semana.
\end{example}

\begin{example}[Liberdade em espaços de funções]\label{ex:b1:vspaces:functionfree}
Em $\mathcal{F}(\R,\R)$, a família $(\eu^{a_1 x}, \dots, \eu^{a_p
x})$ com $a_1 < \dots < a_p$ é \omterm{def:b1:vspaces:free}{livre}: divida uma combinação
nula por $\eu^{a_p x}$ e faça $x \to +\infty$; o último coeficiente morre, e desce-se
em cascata (o \cref{exo:b1:vspaces:8} detalha isto e as suas variantes). A liberdade de
funções demonstra-se \emph{avaliando}: em pontos bem escolhidos, no
infinito, ou depois de derivar.
\end{example}

\begin{example}[Uma relação escondida encolhe um espaço gerado]
\label{ex:b1:vspaces:hiddenrelation}
Em $\mathcal{F}(\R, \R)$, o que é $\operatorname{Vect}(1,\ \cos^2,\ \sin^2)$? A identidade $\cos^2 + \sin^2 = 1$ é uma combinação nula
não trivial
\[
1\cdot\mathbf{1} + (-1)\cos^2 + (-1)\sin^2 = 0 :
\]
a família é linearmente dependente, e o \omterm{def:b1:vspaces:span}{espaço gerado} já é gerado
apenas por $(1, \cos^2)$ ($\sin^2 = 1 - \cos^2$). Essa família menor é \omterm{def:b1:vspaces:free}{livre}: $a + b\cos^2 x = 0$ para todo
$x$ dá, em $x = 0$ e $x = \frac\pi2$: $a + b = 0$ e $a = 0$. Logo o \omterm{def:b1:vspaces:span}{espaço gerado} é
um \emph{plano} dentro do espaço de funções --- e ele contém também
$\cos 2x = 2\cos^2 x - 1$: famílias de funções trigonométricas com ar de independentes
colapsam rotineiramente sob identidades, razão pela qual a liberdade
deve ser \emph{demonstrada}, nunca suposta a partir do tamanho da
lista.
\end{example}

\begin{remark}[Armadilhas comuns]
\label{rem:b1:vspaces:pitfalls}
Quatro armadilhas clássicas.
\emph{Dois a dois não basta}: em $\R^2$, os vetores $(1,0)$,
$(0,1)$, $(1,1)$ são dois a dois não proporcionais, e no entanto
linearmente dependentes --- a liberdade é uma propriedade da família
\emph{inteira}, testada por uma única combinação global, nunca duas
a duas.
\emph{O vetor nulo envenena tudo}: toda família que contém $0$ é
linearmente dependente ($1\cdot 0 = 0$ é uma relação não trivial), por mais
inocentes que sejam os outros vetores.
\emph{União não é soma}: $F \cup G$ quase nunca é um \omterm{def:b1:vspaces:subspace}{subespaço}
(\cref{def:b1:vspaces:subspace}); o menor \omterm{def:b1:vspaces:subspace}{subespaço} que contém ambos é $F + G$, em geral
muito maior do que a união --- em $\R^2$, duas retas distintas têm
por união uma cruz e por soma o plano inteiro.
\emph{Direta exige interseção trivial, não disjunção}: dois
\omterm{def:b1:vspaces:subspace}{subespaços} nunca são disjuntos (ambos contêm $0$); a condição
correta é $F \cap G = \{0\}$, e ela deve ser \emph{demonstrada}, não lida num
desenho --- cf.\ o \cref{ex:b1:vspaces:manysupplements}, em que muitos $G$ diferentes servem.
\emph{A liberdade depende dos escalares}: o par $(1, \iu)$ é \omterm{def:b1:vspaces:free}{livre} em
$\C$ visto como \omterm{def:b1:vspaces:def}{espaço vetorial} sobre $\R$, mas linearmente
dependente em $\C$ visto como \omterm{def:b1:vspaces:def}{espaço vetorial} sobre $\C$ ($\iu\cdot 1 + (-1)\cdot\iu = 0$).
Saiba sempre qual \omterm{def:b1:structures:field}{corpo} está agindo antes de declarar \omterm{def:b1:vspaces:free}{livre} uma
família --- o problema de fim de semana do \cref{ch:b1:findim} transforma
exatamente essa sensibilidade em demonstrações de irracionalidade.
\end{remark}

\begin{remark}[Para onde vai esta linguagem]
\label{rem:b1:vspaces:whereused}
Tudo o que vem depois deste capítulo fala a linguagem instalada
aqui. O \cref{ch:b1:findim} conta vetores de uma \omterm{def:b1:vspaces:free}{base} e transforma ``\omterm{def:b1:vspaces:free}{livre}'' e
``geradora'' em desigualdades sobre um único inteiro, a dimensão.
O \cref{ch:b1:linmaps} estuda as aplicações compatíveis com as duas operações; as
\omterm{def:b1:vspaces:sum}{somas diretas} tornam-se lá projetores. O \cref{ch:b1:matrices} codifica vetores
pelas suas \omterm{prop:b1:vspaces:coordinates}{coordenadas} numa \omterm{def:b1:vspaces:free}{base} --- a \cref{prop:b1:vspaces:coordinates} é a licença para essa
codificação --- e o \cref{ch:b1:euclid} acrescenta comprimentos e ângulos sobre a
estrutura linear. No volume do segundo ano de graduação os mesmos
axiomas, palavra por palavra, valem sobre \omterm{def:b1:structures:field}{corpos} arbitrários e em
dimensão infinita; nada neste capítulo usou a finitude em lugar
algum.
\end{remark}

\begin{remark}[Três fios a seguir ao longo do Livro 3]
\label{rem:b1:vspaces:threads}
Observe três ideias específicas deste capítulo crescerem.
\emph{O princípio da escada}
(\cref{ex:b1:vspaces:staircase}) reaparece como a \omterm{def:b1:vspaces:free}{base} de Newton
no problema de fim de semana deste capítulo, como a \omterm{def:b1:vspaces:free}{base} binomial
$(B_k)$ ali, e como o truque do alternante polinomial no problema de
fim de semana do \cref{ch:b1:det}: um lema, três dividendos sem determinantes.
\emph{A avaliação como teste de liberdade}
(\cref{ex:b1:vspaces:functionfree}) torna-se o isomorfismo de interpolação do \cref{ch:b1:linmaps}, depois o
critério de Vandermonde do \cref{ch:b1:det}, depois o teste de Gram do
\cref{ch:b1:euclid}: o mesmo reflexo, afiado três vezes.
\emph{As \omterm{def:b1:vspaces:sum}{somas diretas}} (\cref{def:b1:vspaces:sum}) tornam-se projetores no \cref{ch:b1:linmaps},
decomposições ortogonais $E =
F \oplus F^\perp$ no \cref{ch:b1:euclid}, e a decomposição
explicado-mais-resíduo dos mínimos quadrados no problema de fim de
semana do \cref{ch:b1:multivar}. Muito pouco deste livro não é, no fundo, uma
destas três ideias vestindo roupa nova.
\end{remark}

\section{Exercícios}

\begin{exercise}[$\star$]\label{exo:b1:vspaces:1}
Quais dos seguintes são \omterm{def:b1:vspaces:subspace}{subespaços}?
\begin{enumerate}
  \item $\{(x, y, z) \in \R^3 : x + 2y - z = 0\}$;
  \item $\{(x, y, z) \in \R^3 : x + 2y - z = 1\}$;
  \item $\{(x, y) \in \R^2 : xy \geq 0\}$;
  \item $\{P \in \R[X] : P(1) = 0\}$;
  \item $\{f \in \mathcal{F}(\R,\R) : f \text{ limitada}\}$.
\end{enumerate}
\end{exercise}

\begin{exercise}[$\star$]\label{exo:b1:vspaces:2}
Em $\R^3$, $(1, 2, 1)$ pertence a
$\operatorname{Vect}\bigl((1,0,1),\, (1,1,0)\bigr)$? E $(2, 1,
1)$? Descreva $\operatorname{Vect}\bigl((1,0,1),(1,1,0)\bigr)$ por uma equação.
\end{exercise}

\begin{exercise}[$\star$]\label{exo:b1:vspaces:3}
Decida a liberdade em $\R^3$: $\;\bigl((1,1,0), (1,0,1),
(0,1,1)\bigr)$; $\;\bigl((1,2,3), (2,4,6)\bigr)$;
$\;\bigl((1,0,0), (1,1,0), (1,1,1), (0,1,1)\bigr)$.
\end{exercise}

\begin{exercise}[$\star$]\label{exo:b1:vspaces:4}
Demonstre que $(1, X - 1, (X-1)^2, (X-1)^3)$ é uma \omterm{def:b1:vspaces:free}{base} de $\R_3[X]$,
e dê as \omterm{prop:b1:vspaces:coordinates}{coordenadas} de $X^3$ nela. \emph{(Taylor em $1$!)}
\end{exercise}

\begin{exercise}[$\star\star$]\label{exo:b1:vspaces:5}
Em $\R^4$, sejam $F = \{(x,y,z,t) : x = y = z\}$ e $G = \{(x,y,z,t)
: x = t = 0\}$. Demonstre que $F \oplus G = \R^4$, e decomponha
$(1,2,3,4)$ em consequência.
\end{exercise}

\begin{exercise}[$\star\star$]\label{exo:b1:vspaces:6}
No espaço das sequências, seja $F$ o \omterm{def:b1:logic:sets}{conjunto} das sequências
convergentes e $G = \operatorname{Vect}(u)$, onde $u_n = (-1)^n$. Demonstre que $F \cap G = \{0\}$. Será $F + G$
o espaço inteiro das sequências?
\end{exercise}

\begin{exercise}[$\star\star$]\label{exo:b1:vspaces:7}
Sejam $F, G, H$ \omterm{def:b1:vspaces:subspace}{subespaços} de $E$. Demonstre que
\[
F \cap (G + (F \cap H)) = (F \cap G) + (F \cap H),
\]
e mostre por um exemplo em $\R^2$ que a distributividade irrestrita
$F \cap (G + H) = (F\cap G) + (F \cap H)$ é falsa.
\end{exercise}

\begin{exercise}[$\star\star\star$]\label{exo:b1:vspaces:8}
Demonstre que as seguintes famílias de $\mathcal{F}(\R, \R)$ são \omterm{def:b1:vspaces:free}{livres}:
\begin{enumerate}
  \item $(\eu^{a_1 x}, \dots, \eu^{a_p x})$ para $a_1 < \dots < a_p$;
  \item $(\cos x, \sin x, \cos 2x, \sin 2x)$;
  \item $(x \mapsto \abs{x - a_1}, \dots, x \mapsto \abs{x - a_p})$
        para $a_i$ distintos \emph{(a \omterm{def:b1:derivative:def}{derivabilidade} falha em
        exatamente um ponto por função)}.
\end{enumerate}
\end{exercise}

\begin{exercise}[$\star\star\star$]\label{exo:b1:vspaces:9}
Seja $E$ um \omterm{def:b1:vspaces:def}{espaço vetorial} sobre $K$ e $F, G, H$ \omterm{def:b1:vspaces:subspace}{subespaços} com
$F + G =
F + H$, $F \cap G = F \cap H$ e $G \subseteq H$. Demonstre que $G = H$.
Dê um contraexemplo sem a hipótese $G \subseteq H$.
\end{exercise}

\begin{exercise}[$\star\star$]\label{exo:b1:vspaces:10}
Em $\R[X]$, seja $\mathcal P$ o \omterm{def:b1:logic:sets}{conjunto} dos \omterm{def:b1:poly:def}{polinômios} \emph{pares}
($P(-X) = P(X)$) e $\mathcal I$ o \omterm{def:b1:logic:sets}{conjunto} dos \emph{ímpares}
($P(-X) = -P(X)$). Demonstre que $\R[X] = \mathcal P \oplus \mathcal
I$, e mostre que $\mathcal P =
\operatorname{Vect}(1, X^2, X^4, \dots)$, isto é, que os
\omterm{def:b1:poly:def}{polinômios} pares são exatamente os \omterm{def:b1:poly:def}{polinômios} em $X^2$.
\end{exercise}

\begin{exercise}[$\star\star$]\label{exo:b1:vspaces:11}
Seja $(x_1, x_2, x_3)$ uma \omterm{def:b1:vspaces:free}{família livre} de um \omterm{def:b1:vspaces:def}{espaço vetorial} real $E$.
Demonstre que $(x_1 + x_2,\ x_2 + x_3,\ x_3 + x_1)$ é \omterm{def:b1:vspaces:free}{livre}. A família análoga de quatro vetores
$(x_1 + x_2,\ x_2 + x_3,\ x_3 +
x_4,\ x_4 + x_1)$ é \omterm{def:b1:vspaces:free}{livre} quando $(x_1, x_2, x_3, x_4)$ o é?
\end{exercise}

\begin{exercise}[$\star\star\star$]\label{exo:b1:vspaces:12}
Seja $E$ um \omterm{def:b1:vspaces:def}{espaço vetorial} sobre $\R$ (ou $\C$) e $F_1, \dots,
F_k$
\omterm{def:b1:vspaces:subspace}{subespaços} \emph{próprios} de $E$ (cada $F_i \neq E$).
\begin{enumerate}
  \item Trate diretamente o caso $k = 2$: se $F_1 \not\subseteq
        F_2$ e $F_2 \not\subseteq F_1$, tome
        $x \in F_1
        \setminus F_2$ e $y \in F_2 \setminus F_1$ e localize $x
        + y$.
  \item Demonstre em geral que $E \neq F_1 \cup \dots \cup F_k$: um
        \omterm{def:b1:vspaces:def}{espaço vetorial} sobre um \omterm{def:b1:structures:field}{corpo} infinito nunca é uma união
        finita de \omterm{def:b1:vspaces:subspace}{subespaços} próprios. \emph{(Tome $k$ mínimo,
        escolha $x \in F_1$ fora dos outros $F_i$, escolha $y \notin F_1$,
        e siga a reta $t \mapsto y + tx$.)}
\end{enumerate}
\end{exercise}

\section{Problema: interpolação, três bases para um espaço}

\begin{problem}\label{pb:b1:vspaces:1}
Fixe $n + 1$ pontos \emph{distintos} $x_0, x_1, \dots, x_n$ de $\R$.
Este problema revisita a \omterm{thm:b1:poly:lagrange}{interpolação de Lagrange}
(o \cref{thm:b1:poly:lagrange}) com os olhos deste capítulo: o
espaço $\R_n[X]$ carrega três \omterm{def:b1:vspaces:free}{bases} naturais --- a de Lagrange,
a de Newton e (para pontos igualmente espaçados) a \omterm{def:b1:vspaces:free}{base} binomial ---
e cada \omterm{def:b1:vspaces:free}{base} torna fácil uma pergunta. O caminho termina num teorema
aritmético genuíno: a caracterização de Pólya dos \omterm{def:b1:poly:def}{polinômios} que
levam $\Z$ em $\Z$.\index{interpolação de Lagrange}
\index{diferenças divididas}\index{polinômio a valores inteiros}

\medskip
\noindent\textbf{Parte I --- A \omterm{def:b1:vspaces:free}{base} de Lagrange.} Para $0 \leq i
\leq n$ ponha
\[
L_i \;=\; \prod_{j \neq i} \frac{X - x_j}{x_i - x_j} \;\in\;
\R_n[X].
\]
\begin{enumerate}
  \item Verifique que $\deg L_i = n$ e que $L_i(x_j) = 1$ se $j =
        i$, e $0$ se $j \neq i$.
  \item Demonstre que a família $(L_0, \dots, L_n)$ é \omterm{def:b1:vspaces:free}{livre}.
  \item Demonstre que para todo $P \in \R_n[X]$,
        \[
        P \;=\; \sum_{i=0}^{n} P(x_i)\, L_i ,
        \]
        e conclua que $(L_0, \dots, L_n)$ é uma \omterm{def:b1:vspaces:free}{base} de
        $\R_n[X]$. \emph{(Considere a diferença dos dois lados
        e conte as suas raízes, \cref{cor:b1:poly:nroots}.)}
  \item Deduza o teorema da interpolação: para quaisquer valores
        $y_0,
        \dots, y_n \in \R$ existe um \emph{único} $P \in
        \R_n[X]$ com $P(x_i) = y_i$ para todo $i$.
        Na \omterm{def:b1:vspaces:free}{base} de Lagrange, quais são as \omterm{prop:b1:vspaces:coordinates}{coordenadas} de um
        \omterm{def:b1:poly:def}{polinômio} $P$?
  \item Demonstre as identidades
        \[
        \sum_{i=0}^{n} L_i = 1
        \qquad\text{e, para } 0 \leq k \leq n,\qquad
        \sum_{i=0}^{n} x_i^{k}\, L_i = X^{k} .
        \]
\end{enumerate}

\medskip
\noindent\textbf{Parte II --- A \omterm{def:b1:vspaces:free}{base} de Newton e as \omterm{pb:b1:vspaces:1}{diferenças
divididas}.} Ponha $N_0 = 1$ e $N_k = (X - x_0)(X - x_1) \cdots
(X - x_{k-1})$ para $1 \leq k \leq n$. Para uma função $f$
definida nos nós, defina as \emph{\omterm{pb:b1:vspaces:1}{diferenças divididas}} por
$f[x_i] = f(x_i)$ e
\[
f[x_i, \dots, x_{i+k}] \;=\;
\frac{f[x_{i+1}, \dots, x_{i+k}] - f[x_i, \dots, x_{i+k-1}]}
{x_{i+k} - x_i} .
\]
\begin{enumerate}[resume]
  \item Demonstre que $(N_0, N_1, \dots, N_n)$ é uma \omterm{def:b1:vspaces:free}{base} de
        $\R_n[X]$.
  \item Calcule $f[x_0, x_1]$ e $f[x_0, x_1, x_2]$ em termos dos
        valores de $f$, e depois calcule todas as \omterm{pb:b1:vspaces:1}{diferenças
        divididas} de $f(x) = x^2$ em três nós arbitrários.
  \item (Lema de Aitken) Seja $R$ interpolando $f$ em $x_0, \dots,
        x_{n-1}$ e $Q$
        interpolando $f$ em $x_1, \dots, x_n$, ambos de grau $\leq n - 1$. Demonstre
        que
        \[
        S \;=\; \frac{(X - x_0)\,Q - (X - x_n)\,R}{x_n - x_0}
        \]
        interpola $f$ em $x_0, x_1, \dots, x_n$.
  \item Deduza, por indução sobre o número de nós, que o
        coeficiente de $X^{k}$ no interpolante de $f$ em $x_0,
        \dots, x_k$
        é exatamente $f[x_0, \dots, x_k]$.
  \item Demonstre a \emph{fórmula de interpolação de Newton}: o
        interpolante de $f$ em $x_0, \dots, x_n$ é
        \[
        P \;=\; \sum_{k=0}^{n} f[x_0, \dots, x_k]\, N_k ,
        \]
        e deduza a fórmula \omterm{def:b1:topology:closed}{fechada}
        \[
        f[x_0, \dots, x_k] \;=\; \sum_{i=0}^{k}
        \frac{f(x_i)}{\prod_{j \neq i,\, j \leq k} (x_i - x_j)} ,
        \]
        que mostra que $f[x_0, \dots, x_k]$ não depende da ordenação dos nós.
\end{enumerate}

\medskip
\noindent\textbf{Parte III --- Nós igualmente espaçados: o operador
de diferenças.} A partir daqui os nós são $0, 1, 2, \dots$ e, para um
\omterm{def:b1:poly:def}{polinômio} $P$, pomos
\[
\Delta P(X) = P(X + 1) - P(X),
\qquad
B_k = \frac{X(X-1)\cdots(X-k+1)}{k!} \quad (B_0 = 1).
\]
\begin{enumerate}[resume]
  \item Mostre que se $\deg P = m \geq 1$ com coeficiente dominante
        $a$, então $\deg \Delta P = m - 1$ com coeficiente dominante
        $m\,a$, e que $\Delta$ mata as constantes.
  \item Mostre que $(B_0, B_1, \dots, B_n)$ é uma \omterm{def:b1:vspaces:free}{base} de
        $\R_n[X]$ e que $\Delta B_k = B_{k-1}$ para $k \geq 1$.
  \item (Fórmula das diferenças progressivas de Newton) Demonstre
        que todo $P
        \in \R_n[X]$ satisfaz
        \[
        P \;=\; \sum_{k=0}^{n} \bigl(\Delta^{k} P\bigr)(0)\, B_k .
        \]
  \item Demonstre que para todo $k \geq 0$,
        \[
        \bigl(\Delta^{k} P\bigr)(0) \;=\;
        \sum_{j=0}^{k} (-1)^{k-j} \binom{k}{j} P(j) .
        \]
  \item Mostre que se $\deg P = n$ com coeficiente dominante $a_n$,
        então $\Delta^{n} P$ é a constante $n!\,a_n$ e
        $\Delta^{n+1} P = 0$.
\end{enumerate}

\medskip
\noindent\textbf{Parte IV --- \omterm{pb:b1:vspaces:1}{Polinômios a valores inteiros}.} Um
\omterm{def:b1:poly:def}{polinômio} $P \in \R[X]$ é \emph{a valores inteiros} quando $P(m) \in
\Z$
para todo $m \in \Z$.
\begin{enumerate}[resume]
  \item Demonstre que cada $B_k$ é a valores inteiros.
        \emph{(Trate $m
        \geq k$, $0 \leq m < k$ e $m < 0$ separadamente; para
        $m =
        -q < 0$, mostre que $B_k(-q) = (-1)^k \binom{q + k - 1}{k}$.)}
  \item Demonstre a \emph{caracterização de Pólya}: $P \in \R_n[X]$
        é a valores inteiros se e somente se as suas \omterm{prop:b1:vspaces:coordinates}{coordenadas} na
        \omterm{def:b1:vspaces:free}{base} $(B_0, \dots, B_n)$ são inteiras.
  \item Deduza: se $P \in \R_n[X]$ assume valores inteiros em $n + 1$
        inteiros \emph{consecutivos} $a, a+1, \dots, a+n$, então $P$
        é a valores inteiros. \emph{(Translade: aplique o estudo a
        $Q(X)
        = P(X + a)$.)}
  \item Deduza da questão 16 que um produto de $k$ inteiros
        consecutivos é sempre divisível por $k!$.
  \item Seja $P = \dfrac{X(X+1)(2X+1)}{6}$. Calcule a sua tabela de Newton em
        $0, 1, 2, 3$, escreva $P$ na \omterm{def:b1:vspaces:free}{base} $(B_k)$, e conclua que $P$ é a
        valores inteiros embora nenhum dos seus coeficientes
        monomiais seja inteiro. Verifique $\Delta P =
        (X+1)^2$ e deduza $P(m) = 1^2 + 2^2 + \dots + m^2$
        para $m \in \N$.
\end{enumerate}

\medskip
\noindent\textbf{Parte V --- Dividendos.}
\begin{enumerate}[resume]
  \item Seja $P \in \R_n[X]$ interpolando os valores $2^i$ em $i =
        0, 1, \dots, n$. Mostre
        que $P = B_0 + B_1 + \dots + B_n$ e que $P(n + 1) = 2^{n+1} - 1$: o ``padrão de duplicação'' quebra
        sempre no ponto imediatamente seguinte.
  \item (Primitiva discreta) Demonstre que para todos os inteiros
        $m
        \geq 1$ e $k \geq 0$,
        \[
        \sum_{j=0}^{m-1} B_k(j) \;=\; B_{k+1}(m),
        \]
        isto é, a identidade do taco de hóquei $\sum_{j=k}^{m-1}
        \binom{j}{k} = \binom{m}{k+1}$.
  \item Expanda $X^2$ e $X^3$ na \omterm{def:b1:vspaces:free}{base} $(B_k)$ e deduza fórmulas
        \omterm{def:b1:topology:closed}{fechadas} para $\sum_{j=0}^{m-1} j^2$ e $\sum_{j=0}^{m-1} j^3$; recupere a identidade de
        Nicômaco $1^3 + \dots + m^3 = (1 + \dots + m)^2$.
  \item Tome $n = 2$ e os nós $0, 1, 2$. Escreva as \omterm{prop:b1:vspaces:coordinates}{coordenadas} de
        $X^2$ nas três \omterm{def:b1:vspaces:free}{bases} deste problema: a \omterm{def:b1:vspaces:free}{base} monomial, a
        \omterm{def:b1:vspaces:free}{base} de Lagrange, a \omterm{def:b1:vspaces:free}{base} de Newton. Confira as três
        respostas contra as questões 4 e 9.
  \item Síntese. Em quatro frases: qual conceito de \omterm{def:b1:vspaces:def}{espaço vetorial}
        torna a questão 4 automática; por que a \omterm{def:b1:vspaces:free}{base} de Newton
        calcula as \omterm{prop:b1:vspaces:coordinates}{coordenadas} \emph{recursivamente} enquanto a
        \omterm{def:b1:vspaces:free}{base} de Lagrange as lê \emph{instantaneamente}; qual
        critério de liberdade as duas \omterm{def:b1:vspaces:free}{bases} partilham; e em que
        sentido preciso o teorema de Pólya diz que a integralidade
        de um \omterm{def:b1:poly:def}{polinômio} é uma propriedade das suas \omterm{prop:b1:vspaces:coordinates}{coordenadas}
        \emph{na \omterm{def:b1:vspaces:free}{base} certa}.
\end{enumerate}
\end{problem}
